\section{Оценивание импульсной характеристики канала}
\label{sec:estimation}

\subsection{Постановка задачи}

Для того чтобы приемное устройство могло правильно демодулировать и декодировать информацию, оно должно знать, как именно пространственный канал исказил переданный сигнал. Для этого в системах связи передаются специальные обучающие последовательности, состоящие из известных приемнику сигналов.

В простейшем случае, когда канал можно считать частотно-неселективным (плоским), а обучающая последовательность состоит всего из одного импульса, принятый сигнал описывается уравнением:
\begin{equation}
    x = h(0) s + z,
    \label{eq:est_simple}
\end{equation}
где $h(0)$ --- комплексный канальный коэффициент (коэффициент передачи прямого сигнала), $s$ --- известный обучающий сигнал, $z$ --- собственный шум приемника. Будем считать дисперсию коэффициента $h(0)$ равной единице, а шум $z$ --- некоррелированным гауссовым процессом с нулевым средним и дисперсией $\sigma_0^2$.

Задача оценивания заключается в нахождении такой оценки $\hat{h}(0)$, которая была бы максимально близка к истинному значению $h(0)$. Для этого используется метод максимального правдоподобия. Функция правдоподобия для (\ref{eq:est_simple}) имеет вид:
\begin{equation}
    p(x) = \frac{1}{\pi\sigma_0^2} \exp\left(-\frac{|x - h(0) s|^2}{\sigma_0^2}\right).
\end{equation}
Максимум этой функции достигается при минимуме показателя экспоненты. Взяв производную по $h(0)$ и приравняв ее к нулю, получим оценку максимального правдоподобия:
\begin{equation}
    \hat{h}(0) = \frac{x s^*}{|s|^2}.
    \label{eq:est_ml_simple}
\end{equation}

Оценка (\ref{eq:est_ml_simple}) является несмещенной, то есть $\langle \hat{h}(0) \rangle = h(0)$. Точность оценивания определяется дисперсией ошибки:
\begin{equation}
    \langle |\hat{h}(0) - h(0)|^2 \rangle = \frac{\sigma_0^2}{|s|^2} = \frac{1}{\rho},
\end{equation}
где $\rho = |s|^2 / \sigma_0^2$ --- отношение сигнал/шум. Из этого выражения видно, что с ростом отношения сигнал/шум дисперсия ошибки оценивания уменьшается.

\subsection{Оценивание с использованием обучающей последовательности}

Чтобы повысить точность оценивания, на практике передают не один, а целую последовательность известных сигналов длиной $L$. В этом случае приемник регистрирует последовательность:
\begin{equation}
    x(k) = h s(k) + z(k), \quad k = 1, 2, \dots, L.
\end{equation}
В векторно-матричной форме это уравнение принимает вид:
\begin{equation}
    \mathbf{X} = h \mathbf{S} + \mathbf{Z},
    \label{eq:est_vector}
\end{equation}
где $\mathbf{X}$, $\mathbf{S}$ и $\mathbf{Z}$ --- векторы-столбцы размерности $L$.

Оценка максимального правдоподобия в этом случае находится по формуле:
\begin{equation}
    \hat{h} = \frac{\mathbf{S}^H \mathbf{X}}{\mathbf{S}^H \mathbf{S}} = \frac{\sum_{k=1}^L x(k) s^*(k)}{\sum_{k=1}^L |s(k)|^2}.
    \label{eq:est_ml_seq}
\end{equation}
Дисперсия этой оценки уменьшается в $L$ раз по сравнению с (\ref{eq:est_ml_simple}):
\begin{equation}
    \langle |\hat{h} - h|^2 \rangle = \left[ \sum_{k=1}^L \frac{|s(k)|^2}{\sigma_0^2} \right]^{-1}.
    \label{eq:est_variance}
\end{equation}
Это означает, что увеличение длины обучающей последовательности повышает точность оценивания канала.

\subsection{Оценивание импульсной характеристики многолучевого канала}

В реальных многолучевых каналах импульсная характеристика содержит не один, а $m+1$ коэффициент: $h(0), h(1), \dots, h(m)$. В этом случае принятый сигнал описывается уравнением свертки:
\begin{equation}
    x(k) = \sum_{n=0}^{m} h(n) s(k - n) + z(k).
    \label{eq:est_multipath}
\end{equation}

Для оценки всех коэффициентов ИХ необходимо использовать более длинную последовательность принятых сигналов. В векторно-матричной форме уравнение (\ref{eq:est_multipath}) записывается аналогично (\ref{eq:est_vector}):
\begin{equation}
    \mathbf{X} = \mathbf{S} \mathbf{H} + \mathbf{Z},
    \label{eq:est_matrix_form}
\end{equation}
где $\mathbf{H} = [h(0), h(1), \dots, h(m)]^T$ --- вектор искомых коэффициентов, а $\mathbf{S}$ --- матрица обучающих векторов, составленная из сдвинутых копий обучающей последовательности.

Оценка максимального правдоподобия (которая в данном случае совпадает с оценкой по методу наименьших квадратов) находится по формуле:
\begin{equation}
    \hat{\mathbf{H}} = (\mathbf{S}^H \mathbf{S})^{-1} \mathbf{S}^H \mathbf{X}.
    \label{eq:est_ls}
\end{equation}
Матрица $(\mathbf{S}^H \mathbf{S})^{-1} \mathbf{S}^H$ называется псевдообратной матрицей к матрице $\mathbf{S}$.

Точность оценивания определяется матрицей ошибок:
\begin{equation}
    \langle \Delta \mathbf{H} \Delta \mathbf{H}^H \rangle = \sigma_0^2 (\mathbf{S}^H \mathbf{S})^{-1} = \sigma_0^2 \mathbf{M}^{-1},
    \label{eq:est_error_matrix}
\end{equation}
где $\mathbf{M} = \mathbf{S}^H \mathbf{S}$ --- матрица Грамма, составленная из скалярных произведений обучающих векторов.

Чтобы оценки разных коэффициентов ИХ были независимыми и имели одинаковую точность, необходимо, чтобы обучающие векторы были ортогональными. В этом случае матрица $\mathbf{M}$ становится диагональной, и матрица ошибок (\ref{eq:est_error_matrix}) также становится диагональной с равными элементами:
\begin{equation}
    \langle |\Delta h(n)|^2 \rangle = \frac{\sigma_0^2}{L |s|^2}.
\end{equation}
Если обучающая последовательность выбрана неудачно (векторы неортогональны), матрица $\mathbf{M}$ становится близкой к сингулярной, что приводит к резкому возрастанию ошибки оценивания. Поэтому выбору обучающей последовательности уделяется большое внимание.

\subsection*{Выводы по разделу}

Таким образом, процедура оценивания канала сводится к решению системы линейных уравнений (\ref{eq:est_matrix_form}) с использованием метода наименьших квадратов. Точность оценивания зависит от отношения сигнал/шум, длины обучающей последовательности и степени ортогональности обучающих векторов. Полученная оценка импульсной характеристики $\hat{\mathbf{H}}$ в дальнейшем используется приемным устройством для эквализации (выравнивания) принятых информационных символов.